% Part: second-order-logic
% Chapter: sol-and-sets
% Section: introduction

\documentclass[../../../include/open-logic-section]{subfiles}

\begin{document}

% SEGMENT OLP-0337-S01

\olfileid{sol}{set}{int}

\olsection{அறிமுகம்}

இரண்டாம்தரத் தருக்கத்தில் பொருட்களத்தின் உட்கணங்கள்மீதும்
சார்புகள்மீதும் அளவையிட முடிவதால், குறைந்தது ஓரளவு கணக் கோட்பாட்டை
அதில் வளர்க்க முடியும் என்று எதிர்பார்க்கலாம். “வளர்த்தல்” என்பதால்,
கணக் கோட்பாட்டுப் பண்புகளையும் கூற்றுகளையும் இரண்டாம்தரத் தருக்கத்தில்
வெளிப்படுத்த முடியும் என்றும், கணங்களுக்கான தனிப்பட்ட தருக்கமல்லாத
சொற்களஞ்சியம் (எடுத்துக்காட்டாக, கணக் கோட்பாட்டின் உறுப்புறவு
!!{predicate}) எதுவுமின்றி அதைச் செய்ய முடியும் என்றும் குறிக்கிறோம்.
எடுத்துக்காட்டாக, கணங்களின் ஒன்றிப்புகளையும் வெட்டுகளையும் உட்கண
உறவையும் வரையறுக்கலாம்; கணங்களின் அளவுகளையும் ஒப்பிட்டு, கான்டரின்
தேற்றம் போன்ற முடிவுகளையும் கூறலாம்.

\end{document}
